\bookchapter{Kapitel 33: Zukunftsausblick – Die naechste Etappe des Smartphones}

\chapterintro{Das Handy wird nicht ueber Nacht verschwinden, aber seine dominante Stellung koennte von einer neuen Geraetegeneration aufgeteilt werden. In diesem Kapitel besprechen wir drei Dinge: Wie sehen die Geraete der naechsten Generation aus, wie wird das Handy als neues Fundament der Identitaet »unsichtbar«, und welche ethischen Fragen muessen wir im Voraus bedenken.}

\booksection{AR/VR und Spatial Computing: Der »zweite Desktop« neben dem Handy}

Der Wettbewerb um die naechste Geraetegeneration hat begonnen. Die wahrscheinlichsten Kandidaten, die Handy-Nutzungszeit aufzuteilen, sind AR/VR- und Spatial-Computing-Geraete:

\begin{itemize}
  \item \textbf{Vision Pro Folgegeneration}: Der Apple Vision Pro markiert die offizielle Ankunft des Konzepts "Spatial Computing". Aber Preis (ab 29.999 Yuan) und Gewicht der ersten Generation sind noch Hürden. Für 2026–2027 wird eine leichtere, günstigere Version erwartet, die Spatial Computing vom "Entwicklerspielzeug" zum "frühen Massenmarkt" bringt. Dann könnte das iPhone zum Rechen- und Daten-Hub der Vision-Brille werden.
  \item \textbf{Vom "Geräte-Zentrum" zum "Identitäts-Zentrum"}: Ein "zweiter Desktop" vor Ihren Augen – digitale Informationen allgegenwärtig wie Luft. Die Smart-Glasses-Kooperation von Meta und Ray-Ban hat bereits Millionen verkauft. Auch chinesische Hersteller wie Rokid, XREAL und INMO bringen AR-Brillen für Endkunden. 2025–2026 können AR-Brillen schon Echtzeitübersetzung, HUD-Navigation und AI-Assistent-Dialog. Aber sie sind durch Lichtwellenleiter-Displays und Akkulaufzeit begrenzt – 3–5 Jahre bis sie das Handy "ersetzen" können.
  \item \textbf{Smart-Ringe und Wearables}: Samsung Galaxy Ring, Oura Ring und andere Smart-Ringe werden beliebt – als "Körper-Sensor" des Handys für Gesundheitsdaten. Die Rolle des Handys wandelt sich vom "Universal-Terminal" zu einem "Hub + mehrere leichte Zubehörteile".
\end{itemize}

\booksection{Brain-Computer Interface: Gedanken als Interaktionsmedium}

Gedanken selbst werden zum Interaktionsmedium, aber dieser Weg ist laenger als gedacht:

\begin{itemize}
  \item \textbf{Neuralink – neuester Stand}: Musks Neuralink hat 2024 die erste menschliche Implantation durchgeführt. Der erste Proband Noland Arbaugh steuerte per Gedanken eine Computer-Maus und spielte Schach. Aktuell richten sich implantierbare BCIs (Brain-Computer-Interfaces) hauptsächlich auf Hilfen für Gelähmte, mit dem Ziel, ihnen die "Bedienung" der digitalen Welt zurückzugeben.
  \item \textbf{Nicht-invasive BCI}: Mehrere Firmen entwickeln nicht-invasive Kopfbedeckungen (wie Kernel, OpenBCI), die ohne Operation Gehirnsignale lesen, aber viel ungenauer sind als invasive. Noch weit entfernt vom "Gedanken-Tippen".
  \item \textbf{Verbraucher-Ausblick}: BCI ist 10–20 Jahre von der Handy-Interaktion für Endverbraucher entfernt. Kurzfristig könnte es zuerst bei Barrierefreiheit und Aufmerksamkeits-Monitoring eingesetzt werden. Ethik- und Privatsphäre-Fragen ("Wem gehören meine Gehirnwellen?") werden vor der technischen Reife zum Streitpunkt.
\end{itemize}

\booksection{Das »Verschwinden« des Handys und das neue Fundament der Identitaet}

Eine andere kuenftige Entwicklungsrichtung des Handys koennte das »Verschwinden« sein – als allgegenwaertiger »unsichtbarer Helfer« wird es zum einzigen Fundament deiner digitalen Identitaet:

\begin{itemize}
  \item \textbf{Viertens. Verbunden bleiben und trotzdem unabhängig – deine Philosophie des digitalen Lebens}: Sofortige "Du bist Du"-Authentifizierung – bei keinem Dienst mehr manuell Codes eingeben.
  \item \textbf{Lass die Automatisierung dein treuer "Copilot" sein.}: Beim Aktivieren eines neuen Geräts wird das digitale Erbe voreingestellt – "Nachlass" als Werkseinstellung.
  \item \textbf{Führe regelmäßig "digitales Ausmisten" bzw. "soziales Fasten" durch.}: Ihre persönliche Cloud-Datentresor ist das Zentrum; Hardware ist nur der Zugang. Beim Gerätewechsel bleiben Identität und Daten bei Ihnen.
\end{itemize}

\booksection{Ethische Fragen: Sind wir noch die Herren unserer Handys?}

Schreibe für dich und dein Smartphone eine kurze “Nutzungsvereinbarung”, die eure Beziehung definiert. Zum Beispiel:

\begin{itemize}
  \item \textbf{AI-Ethik}: Wenn On-Device-AI Ihre Vorlieben besser kennt als Sie selbst: Ist Ihre "Wahl" noch frei oder gelenkt? Wenn ein AI-Agent in Ihrem Namen bestellt und antwortet: Wo liegen Verantwortung und Grenzen?
  \item \textbf{Datenhoheit}: Wem gehören Ihre Gesundheits-, Gehirnwellen- und Biometrie-Daten – Ihnen, dem Gerätehersteller oder dem Cloud-Anbieter? Das chinesische Gesetz zum Schutz persönlicher Daten und die EU-DSGVO sind erst der Anfang. "Datenübertragbarkeit" und "Recht auf Vergessenwerden" sind noch lange nicht umgesetzt.
  \item \textbf{Digitale Kluft}: AI-Handys, faltbare Displays und Satellitenkommunikation vergrößern die Lücke zwischen "kann sich leisten" und "kann sich nicht leisten". Der gerechte Zugang zur digitalen Infrastruktur wird zu einem neuen gesellschaftlichen Thema – so entscheidend wie vor 20 Jahren die Frage "Kann man ins Internet?".
  \item §1 Meine Aufmerksamkeitszeit gehört mir. Du sollst bei der Arbeit / in der Zeit mit der Familie von selbst in den "Nicht-stören"-Zustand wechseln.
  \item §2 Meine persönlichen Daten sind mein Eigentum. Deine Hauptaufgabe ist sichere Verschlüsselung und komfortable Synchronisation.
\end{itemize}

\booksection{Verbunden bleiben und trotzdem unabhaengig – deine Philosophie des digitalen Lebens}

Egal wie sich die Technik entwickelt – diese Prinzipien digitalen Lebens werden auch in 10 Jahren noch gelten:

\begin{enumerate}
  \item Fünftens. Anhang
  \item Geschafft! Nach dem "langen Marsch" durch die vorangegangenen über zwanzig Kapitel bist du vom "Smartphone-Anfänger" erfolgreich aufgestiegen und beherrschst die Kernkompetenzen vom Einstieg bis zur Meisterschaft. Dieser Anhang ist wie ein Taschenhandbuch und ein Ressourcenschatz für dich.
  \item Anhang A: "Babysitter"-Ablauf für ein neues Smartphone
  \item 1. Auspacken und Prüfen in drei Schritten
  \item Äußerliches prüfen
\end{enumerate}

\begin{tcbpractice}{: Gehäuse und Bildschirm auf Kratzer und Dellen kontrollieren.}
Schreibe fuer dich und dein Handy (und kuenftige AR-Brillen, AI-Assistenten) ein einfaches »Nutzungsabkommen«, das eure Beziehung definiert. Zum Beispiel:

\begin{itemize}
  \item : IMEI/Seriennummer auf der Verpackung mit der auf der Rechnung vergleichen.
  \item Funktionen testen
  \item : Einen Anruf tätigen, um Hörer und Mikrofon zu prüfen; ein Foto machen, um die Kamera zu testen.
\end{itemize}

\end{tcbpractice}

\begin{tcbsummary}{Kapitelzusammenfassung}
\begin{itemize}
  \item Vision Pro und AR-Brillen sind neue Computing-Zugaenge neben dem Handy, koennen es aber kurzfristig nicht ersetzen
  \item BCI ist noch 10–20 Jahre von der Consumer-Reife entfernt; ethische Fragen reifen langsamer als die Technik
  \item Das Handy entwickelt sich in Richtung »Verschwinden«: vom Geraete-Zentrum zum Identitaets-Zentrum
  \item AI-Ethik, Datenhoheit und digitale Kluft sind die Kernthemen der naechsten 10 Jahre
  \item Verbunden bleiben und unabhaengig bleiben – die unwandelbare Philosophie des digitalen Lebens
\end{itemize}

\end{tcbsummary}
